\documentclass[addpoints]{exam}
% \documentclass[addpoints,12pt]{exam}

\usepackage[slovene]{babel}
\usepackage{amsfonts}
\usepackage{amssymb}
\usepackage{amsmath}
\usepackage{float}
\usepackage{graphicx}
\usepackage{enumitem}
\usepackage{color}
\usepackage{eurosym}

\usepackage{pgfpages}
% \usepackage{tikz}
\usepackage{wrapfig}
\usepackage{pgfkeys}
\usepackage{pgfplots}
\usepackage{xcolor}
\usepackage{tkz-euclide}
\usepackage{xfp}




%Komentar naslednje vrstice glede na verzijo testa -- rešitve ali prostor za odgovore:
% \printanswers

% \linespread{2}


\pointsinrightmargin
\marginbonuspointname{}


\renewcommand{\solutiontitle}{\noindent\textbf{Rešitve in točkovanje:}\par\noindent}

\colorgrids
\definecolor{GridColor}{gray}{0.7}
\setlength{\gridsize}{7mm}

\extrawidth{5mm}
\setlength{\rightpointsmargin}{1.5cm}

\begin{document}

\runningfooter{}{}{}

\begin{center}
    \fbox{\fbox{\parbox{15cm}{\centering
    Popravljanje in pridobivanje ocen -- drugo pisno ocenjevanje znanja \\ 2. letnik -- test I \\ 18. junnij 2026 \\ (Kotne funkcije; Vektorji)}}}
\end{center}

\vspace{0.1in}
\makebox[\textwidth]{Ime in priimek, oddelek:\enspace\hrulefill}


\begin{center}
    \hqword{Naloga}
    \hpword{Možne točke}
    \chbpword{Dodatne točke}
    \hsword{Dosežene točke}
    \htword{Skupaj}  
    % \combinedpointtable[h][questions]
    \gradetable[h][questions]

\end{center}

\vspace{0.1in}


\begin{questions}

    % 1. vprašanje

    \question[4]
    Poenostavite izraz. 
    $$ \left(1 + \tan^2x\right)\cdot\cos^2x-\sin^2x(1+\cot^2x)$$

    \begin{solution}[6.5cm]
       
    \end{solution}


    % 2. vprašanje

    \question[4]
    Izračunajte natančno vrednost izraza. 
    $$ \dfrac{\sin^2{135^\circ}+\cos^2{150^\circ}}{\tan\left(-330^\circ\right)}$$

    \begin{solution}[6cm]
       
    \end{solution}


\newpage
    % 3. vprašanje

    \question
    V kocki $ABCDEFGH$ je točka $M$ razpolovišče roba $FG$, točka $N$ razdeli rob $CD$ v razmerju \\ $|CN|:|DN|=1:2$.
    Z baznimi vektorji $\overrightarrow{a}=\overrightarrow{AB}, \overrightarrow{b}=\overrightarrow{AD}$ in $\overrightarrow{c}=\overrightarrow{AE}$ izrazite vektorje: 
    \begin{parts}
        \part[2]
        $\overrightarrow{NM}$, 

    \begin{solution}[5cm]

    \end{solution}
        
        \part[3] 
        $\overrightarrow{MB}-\overrightarrow{HB}+\overrightarrow{HA}+\overrightarrow{DH}$.
    \end{parts}

    \begin{solution}[5cm]

    \end{solution}

    

    
    % 4. vprašanje

    \question[3]
    V trikotniku $\triangle ABC$ z oglišči $A(4,5,-6)$, $B(1,7,0)$ in $C(3,7,-4)$ izračunajte skalarni produkt $\overrightarrow{AB}\cdot\overrightarrow{AC}$.
    % \begin{parts}

    %     \part[2]
    %     dolžino težiščnice na stranico $AC=b$, 

    % \begin{solution}[6cm]

    % \end{solution}

    %     \part[3] 
    %     koordinate težišča trikotnika $\triangle ABC$.

    % \begin{solution}[2cm]

    % \end{solution}        
    
    % \part[5]
    %     skalarni produkt $\overrightarrow{AB}\cdot\overrightarrow{AC}$ in kot $\angle BAC=\alpha$, 

    %     \begin{solution}[5cm]

    % \end{solution}




    % \end{parts}







    \newpage
    % 5. vprašanje

    \question
    Daljica $AB$ je določena s krajiščema $A(5,7,-2)$ in $B(1,11,5)$.
    
    \begin{parts}
        \part[3]
        Določite koordinate točke $T$, ki leži na nosilki daljice $AB$, da bo veljalo $|AB|:|BT|=1:2$.
       
           \begin{solution}[5cm]

    \end{solution}

        \part[3]    
        Določite enotski vektor v smeri vektorja $\overrightarrow{BA}$.

    \begin{solution}[5cm]

    \end{solution}

        \part[3]
        Določite $x$, da bo vektor $\overrightarrow{v}=(x,-1,-1)$ pravokoten na krajevni vektor točke $A$.

            \begin{solution}[5cm]

    \end{solution}

        \part[3]
        Določite $m$ in $n$, da bo vektor $\overrightarrow{u}=(m+1,n,10)$ kolinearen s krajevnim vektorjem točke $B$.
    \end{parts}

\begin{solution}[4cm]

\end{solution}


\newpage

    % 6. vprašanje



    \question[4]
    Dolžina vektorja $\overrightarrow{a}$ je $4$, dolžina vektorja $\overrightarrow{b}$ je $3$, 
    kot med njima pa meri $30^\circ$. 
    Izračunajte skalarni produkt vektorjev $\overrightarrow{a}$ in $\frac{1}{2}\overrightarrow{a}+\frac{4}{3}\overrightarrow{b}$.
    % \begin{parts}
    %     \part[4]
    %     Izračunajte dolžino vektorja $\overrightarrow{c}=\frac{1}{2}\overrightarrow{a}+\frac{4}{3}\overrightarrow{b}$.
        
    %         \begin{solution}[5cm]

    % \end{solution}

    %     \part[3]
    %     Izračunajte skalarni produkt vektorjev $\overrightarrow{a}$ in $\frac{1}{2}\overrightarrow{a}+\frac{4}{3}\overrightarrow{b}$.

    % \end{parts}


    \begin{solution}[10cm]

    \end{solution}





    \question[3]
    Dana sta vektorja $\overrightarrow{a}=(\sin x, -\cos x, 1)$ in $\overrightarrow{b}=(-\sin x, \cos x, 1)$.
    Izračunajte dolžino obeh vektorjev in pokažite, da sta pravokotna za vsak $x\in\mathbb{R}$.

    \begin{solution}[1cm]

    \end{solution}


\end{questions}



\end{document}